\newcommand{\MercurySymbol}{\text{\mdseries\AstroSymbols ☿}}
\bigskip\Needspace{10\baselineskip}
\pdfbookmark[1]{Merkurdurchgang vor der Sonne am 3. November 1802}{mercur18021103}
\begin{center}
MERKURDURCHGANG VOR DER SONNE AM 3. NOVEMBER 1802.\par
\medskip\hrule\medskip
Monatliche Correspondenz. Band VI. S.\ 567—575. Dezember 1802.\par
\medskip\hrule
\end{center}\medskip
—\quad —\quad —

In \textit{Braunschweig} beobachtete Dr.\ \textsc{Gauß} beim Austritt die \textit{innere} Berührung um $0^{\mathrm u}\ 39^{\mathrm m}\ 16^{\mathrm s}$ und die \textit{äußere} um $0^{\mathrm u}\ 40^{\mathrm m}\ 48^{\mathrm s}$ wahrer Zeit. [Handschrift. Austritt des $\MercurySymbol$ in Braunschweig am 9. Nov.\ beobachtet. $0^{\mathrm u}\ 23^{\mathrm m}\ 17^{\mathrm s}$ innere Berührung, $0^{\mathrm u}\ 24^{\mathrm m}\ 49^{\mathrm s}$ äußere Berührung, mittlere Zeit.] Er verwendete dazu einen von \textsc{Baumann} gefertigten zweifüßigen Achromaten, in dem Merkur und die schöne Gruppe von Sonnenflecken, an der er vorbeizog, sehr gut zu sehen waren. Leider konnte er die Zeit nicht so genau bestimmen, wie er es sich gewünscht hätte. Diese Schwierigkeit teilt Dr.\ \textsc{Gauß}
\SourceLeaf{237}%
zu dieser Jahreszeit mit mehreren Astronomen, die kein Passageninstrument haben. So musste er sich zur Zeitbestimmung mit einigen einzelnen Sonnenhöhen begnügen, die er nach 1 Uhr mit einem unbedeckten Ölhorizont gemessen hatte. Aus diesem Grund und weil dies seine erste Beobachtung dieser Art ist, bittet er um Nachsicht, falls seine Beobachtung mit anderen nicht genau übereinstimmen sollte. Außer der Beobachtung des Austritts versuchte er auch, das Fernrohr als Kreismikrometer zu verwenden, und bestimmte mehrere Unterschiede der Rektaszension zwischen Merkur und dem Mittelpunkt der Sonne. Er selbst hält diese Werte allerdings für wenig genau. Immerhin lassen sie sich durch Rechnung überprüfen.
\begin{center}
\begin{tabular}{r|r}
Wahre Zeit&\shortstack{RA-Differenz\\$\odot-\MercurySymbol$}\\\hline
$23^{\mathrm u}\ 26^{\mathrm m}\ 6^{\mathrm s}$&$7'\ 51''$\\
$23\ 42\ 30$&$9\ 29$\\
$23\ 50\ 28$&$10\ 3$\\
$0\ 3\ 25$&$11\ 37$\\
$0\ 25\ 14$&$13\ 41$
\end{tabular}
\end{center}
—\quad —\quad —
\begin{flushright}[\textsc{von Zach.}]\end{flushright}

